%!TEX root = ../main.tex
\chapter{Guía y registro de las entrevistas a usuarios referentes}\label{apx:entrevistas}

\section{Objetivo y modalidad}

Las entrevistas tuvieron como objetivo relevar la percepción de usuarios clave sobre el proceso de gestión de crédito y facturación antes y después de la implementación de la solución, y validar cualitativamente los tiempos estimados en el Capítulo 5. Se utilizó un formato semiestructurado, con una guía de preguntas organizada en cuatro ejes que se adaptó al rol y a la operación de cada entrevistado. Al inicio de cada entrevista se explicó el objetivo del Proyecto Final, se solicitó autorización para tomar notas y se aclaró que en el documento no se consignaría el nombre del entrevistado, sino únicamente su rol.

\section{Guía de la entrevista}

Eje 1: situación previa

\begin{enumerate}
\def\labelenumi{\arabic{enumi}.}
\item
  ¿Cómo era el proceso de carga del crédito de los concesionarios antes del sistema?
\item
  ¿Qué parte del proceso resultaba más tediosa o problemática?
\end{enumerate}

Eje 2: tiempos

\begin{enumerate}
\def\labelenumi{\arabic{enumi}.}
\setcounter{enumi}{2}
\item
  ¿Cuánto tiempo llevaba cargar el saldo de un concesionario?
\item
  ¿Cuánto tiempo llevaba armar el archivo de Excel para Santander por concesionario?
\item
  En un cierre de mes, ¿cuánto tiempo se destinaba a estas tareas?
\item
  ¿Cuánto tiempo lleva hoy esa misma tarea con el sistema?
\end{enumerate}

Eje 3: errores y confiabilidad

\begin{enumerate}
\def\labelenumi{\arabic{enumi}.}
\setcounter{enumi}{6}
\item
  ¿Con qué frecuencia se producían errores de carga? ¿Recuerda algún caso concreto?
\item
  Cuando ocurría un error, ¿qué había que hacer para corregirlo?
\item
  ¿Considera que el saldo que muestra hoy la pantalla CR08 es correcto?
\end{enumerate}

Eje 4: experiencia de uso

\begin{enumerate}
\def\labelenumi{\arabic{enumi}.}
\setcounter{enumi}{9}
\item
  ¿Qué opina de visualizar los saldos de livianos y pesados en una misma pantalla CR08?
\item
  ¿Le resulta útil el correo de despacho? ¿Cambió la forma en que sigue la facturación?
\item
  ¿Qué mejoraría o qué le gustaría que el sistema hiciera y hoy no hace?
\end{enumerate}

La Tabla~\ref{tab:20} resume los datos de cada entrevista.

\section{Datos de las entrevistas}

\begin{longtable}[]{@{}
  >{\raggedright\arraybackslash}p{(\columnwidth - 12\tabcolsep) * \real{0.1561}}
  >{\raggedright\arraybackslash}p{(\columnwidth - 12\tabcolsep) * \real{0.1676}}
  >{\raggedright\arraybackslash}p{(\columnwidth - 12\tabcolsep) * \real{0.1321}}
  >{\raggedright\arraybackslash}p{(\columnwidth - 12\tabcolsep) * \real{0.1335}}
  >{\raggedright\arraybackslash}p{(\columnwidth - 12\tabcolsep) * \real{0.1209}}
  >{\raggedright\arraybackslash}p{(\columnwidth - 12\tabcolsep) * \real{0.1208}}
  >{\raggedright\arraybackslash}p{(\columnwidth - 12\tabcolsep) * \real{0.1690}}@{}}
\caption{Datos de las entrevistas}\label{tab:20}\\
\toprule\noalign{}
\begin{minipage}[b]{\linewidth}\raggedright
Entrevistado
\end{minipage} & \begin{minipage}[b]{\linewidth}\raggedright
Rol
\end{minipage} & \begin{minipage}[b]{\linewidth}\raggedright
Operación
\end{minipage} & \begin{minipage}[b]{\linewidth}\raggedright
Temas abordados
\end{minipage} & \begin{minipage}[b]{\linewidth}\raggedright
Fecha
\end{minipage} & \begin{minipage}[b]{\linewidth}\raggedright
Duración
\end{minipage} & \begin{minipage}[b]{\linewidth}\raggedright
Modalidad
\end{minipage} \\
\midrule\noalign{}
\endfirsthead
\caption[]{Datos de las entrevistas (continuación)}\\
\toprule\noalign{}
\begin{minipage}[b]{\linewidth}\raggedright
Entrevistado
\end{minipage} & \begin{minipage}[b]{\linewidth}\raggedright
Rol
\end{minipage} & \begin{minipage}[b]{\linewidth}\raggedright
Operación
\end{minipage} & \begin{minipage}[b]{\linewidth}\raggedright
Temas abordados
\end{minipage} & \begin{minipage}[b]{\linewidth}\raggedright
Fecha
\end{minipage} & \begin{minipage}[b]{\linewidth}\raggedright
Duración
\end{minipage} & \begin{minipage}[b]{\linewidth}\raggedright
Modalidad
\end{minipage} \\
\midrule\noalign{}
\endhead
\bottomrule\noalign{}
\endlastfoot
E1 & Analista contable, área de Contabilidad General (Accounting Analyst -- General Accounting) & Chile y Perú & Ejes 1 a 4 completos & 6/10/2026 & 45 min aprox. & Presencial \\
E2 & Integrante del área de Gerencia de Ventas (General Sales -- Management) & Perú & Carga de DUA mediante archivo, planilla de Excel para Santander y cantidad de unidades por lote & 8/10/2026 & 15 a 20 min aprox. & Videollamada (Microsoft Teams) \\
\end{longtable}

La Tabla~\ref{tab:21} presenta una síntesis de las respuestas de cada entrevistado, organizada según los ejes de la guía.

\section{Registro de las respuestas}

\begin{longtable}[]{@{}
  >{\raggedright\arraybackslash}p{(\columnwidth - 4\tabcolsep) * \real{0.3333}}
  >{\raggedright\arraybackslash}p{(\columnwidth - 4\tabcolsep) * \real{0.3334}}
  >{\raggedright\arraybackslash}p{(\columnwidth - 4\tabcolsep) * \real{0.3334}}@{}}
\caption{Registro de las respuestas por eje}\label{tab:21}\\
\toprule\noalign{}
\begin{minipage}[b]{\linewidth}\raggedright
Eje
\end{minipage} & \begin{minipage}[b]{\linewidth}\raggedright
E1 (Contabilidad General)
\end{minipage} & \begin{minipage}[b]{\linewidth}\raggedright
E2 (Gerencia de Ventas, Perú)
\end{minipage} \\
\midrule\noalign{}
\endfirsthead
\caption[]{Registro de las respuestas por eje (continuación)}\\
\toprule\noalign{}
\begin{minipage}[b]{\linewidth}\raggedright
Eje
\end{minipage} & \begin{minipage}[b]{\linewidth}\raggedright
E1 (Contabilidad General)
\end{minipage} & \begin{minipage}[b]{\linewidth}\raggedright
E2 (Gerencia de Ventas, Perú)
\end{minipage} \\
\midrule\noalign{}
\endhead
\bottomrule\noalign{}
\endlastfoot
Situación previa & En Chile, los analistas enviaban por correo las líneas disponibles como imagen, y cada saldo se transcribía a mano («se debía cargar cada línea de crédito a ojo»). Los concesionarios con livianos y pesados requerían carga en dos pantallas. En Perú el proceso era más simple: una sola línea, pocos concesionarios y datos tomados del home banking o recibidos por correo. & No se abordó \\
Tiempos & No precisó un tiempo por concesionario, porque varía según el momento del mes y la cantidad de unidades. El cierre de mes concentra la mayor carga, con picos a principio o a mitad de mes. «Se llega a perder horas semanales». Con la solución, la actualización del saldo se realiza «en segundos». & El armado del Excel dependía de la cantidad de unidades y podía superar los diez minutos por concesionario: había que cargar el precio, la fecha de factura y los datos aduaneros de cada vehículo, en un archivo distinto por concesionario. \\
Errores y confiabilidad & Errores de carga frecuentes en Chile, donde los saldos superan las diez cifras: «Muy común olvidarse de poner un 0». Interrumpían la facturación por saldo insuficiente y obligaban a investigar y recalcular. Eran más frecuentes en los cierres. Considera confiable el saldo actual, aunque el equipo lo sigue verificando contra el home banking. & No se abordó \\
Experiencia de uso & La unificación en CR08 es «un gran cambio». Valora poder modificar el saldo a mano si falla la comunicación con el banco. El correo de despacho es «muy útil para saber cuándo finalizó el proceso de facturación», y su demora funciona como alerta. También lo usan los despachantes de Chile y Forum. No identificó mejoras adicionales. & Propuso eliminar el límite de unidades por solicitud a Santander y que el desembolso no requiera aprobar cada planilla en el home banking. \\
\end{longtable}
